\section{System Architecture}

The system has 15 containers on one Docker bridge network called
\texttt{mic-net}. Containers find each other by service name through Docker's
internal DNS, so \texttt{gen-http2} can simply call
\texttt{http://target-http2:8080} without anyone configuring an IP address
anywhere.

Ten of those containers carry the core pipeline described below; the
remaining five are the Mosquitto broker and four \texttt{tshark}-based
analyzer sidecars (one per target's network namespace), introduced at the
end of this section. The ten core containers fall into four groups. The \textbf{controller} (FastAPI,
port 8000) is the brain: it reads the YAML profile, drives the four generators
through their REST APIs, and runs the Adaptive Control loop. This loop is a
real feedback mechanism, not just a timer working through YAML phases. Every
few seconds the controller checks each generator's current error rate against
a threshold and, if it is crossed, scales that generator's rate down by a
bounded multiplier on its own, no human involved. Section III-E gives the full
description. The four \textbf{generators} (\texttt{gen-http2}, \texttt{gen-quic},
\texttt{gen-mqtt}, \texttt{gen-tcpudp}) each keep their own state and expose
the same four endpoints: start, stop, patch the config, get the status. They
are built on \texttt{httpx} (HTTP/2), \texttt{aioquic} \cite{aioquic} (QUIC),
\texttt{paho-mqtt} \cite{pahomqtt} (MQTT) and plain Python sockets (raw
TCP/UDP). The three \textbf{target services} (\texttt{target-http2},
\texttt{target-quic}, \texttt{target-tcpudp}) receive traffic alongside the
separately-counted Mosquitto broker \cite{mosquitto}. \texttt{target-http2} runs on Hypercorn
\cite{hypercorn}, the only ASGI server in the project that supports HTTP/2
natively. The \textbf{metrics collector} (Flask, port 9090) aggregates
everything into one \texttt{/metrics} response, and the \textbf{dashboard}
(a single static HTML file served over port 3000) polls the controller every
two seconds and renders the live state.

Reporting to the metrics collector is not symmetric across these
containers, and it is worth being precise about who reports and why. All four
generators post their own send-side counters (packets sent, error count,
current rate) to the collector, since the controller's Adaptive Control loop
reads those numbers. \texttt{target-http2} and \texttt{target-quic} also post
what they received, because both run custom code that can be instrumented
freely. Mosquitto cannot do this: it runs the stock
\texttt{eclipse-mosquitto:2} image with no custom instrumentation hook, so it
stays a black box. Every MQTT number in this report, throughput, latency,
QoS-level counts, comes only from \texttt{gen-mqtt} on the sending side, never
from the broker itself. \texttt{target-tcpudp} stays equally silent, but for
a different reason: it is meant to disappear cleanly when stopped, and the
resulting TCP RST is exactly what Section VI needs to see in a Wireshark
capture. A reporting connection on top of that would just be one more thing
that breaks when the container goes down, so it was left out on purpose.

Fig.~\ref{fig:architecture} shows how a request actually travels through the
system. The dashboard never talks to a generator directly; every change goes
through the controller, which is what makes it possible to log every
configuration change with a timestamp, as the assignment requires.

\begin{figure*}[t]
\centering
\resizebox{\textwidth}{!}{%
\begin{tikzpicture}[
  box/.style={draw, rounded corners, minimum width=2.6cm, minimum height=0.7cm, align=center, font=\footnotesize},
  arr/.style={-{Latex[length=2mm]}, thick},
]
\node[box] (dash) at (5.2,6.6) {Dashboard\\:3000};
\node[box] (ctrl) at (5.2,5.2) {Controller\\:8000};

\node[box] (h2)     at (0,3.4) {gen-http2};
\node[box] (quic)   at (3.4,3.4) {gen-quic};
\node[box] (mqtt)   at (6.8,3.4) {gen-mqtt};
\node[box] (tcpudp) at (10.2,3.4) {gen-tcpudp};

\node[box] (th2)    at (0,1.6) {target-http2};
\node[box] (tquic)  at (3.4,1.6) {target-quic};
\node[box] (broker) at (6.8,1.6) {mosquitto};
\node[box] (ttcp)   at (10.2,1.6) {target-tcpudp};

\node[box, minimum width=15cm, minimum height=0.8cm] (metrics) at (5.5,-0.5) {Metrics Collector \;\textbar\; :9090};

\draw[arr] (dash) -- (ctrl) node[midway, right, font=\scriptsize] {GET /status};
\draw[arr] (ctrl) -- (h2);
\draw[arr] (ctrl) -- (quic);
\draw[arr] (ctrl) -- (mqtt);
\draw[arr] (ctrl) -- (tcpudp);
\draw[arr] (h2) -- (th2);
\draw[arr] (quic) -- (tquic);
\draw[arr] (mqtt) -- (broker);
\draw[arr] (tcpudp) -- (ttcp);

\draw[arr, dashed] (th2) -- (0,-0.1);
\draw[arr, dashed] (tquic) -- (3.4,-0.1);

\draw[arr, dashed] (h2.south)     -- (0,2.9)  -- (1.7,2.9) -- (1.7,-0.1);
\draw[arr, dashed] (quic.south)   -- (3.4,2.9) -- (5.1,2.9) -- (5.1,-0.1);
\draw[arr, dashed] (mqtt.south)   -- (6.8,2.9) -- (8.5,2.9) -- (8.5,-0.1);
\draw[arr, dashed] (tcpudp.south) -- (10.2,2.9) -- (12.3,2.9) -- (12.3,-0.1);

\draw[arr, dashed] (ctrl.east) -- (13.8,5.2) -- (13.8,-0.5)
  node[pos=0.4, right, font=\scriptsize] {GET /metrics} -- (metrics.east);
\end{tikzpicture}%
}
\caption{Component diagram of the single-machine deployment. Solid arrows are
REST calls or raw protocol traffic. Dashed arrows show metrics reporting: all
four generators and the two custom-built targets post their counters up to
the collector; the long line on the right is the controller reading the
aggregate back out for the dashboard and for Adaptive Control. Mosquitto and
\texttt{target-tcpudp} have no dashed line on purpose, as explained in the
text.}
\label{fig:architecture}
\end{figure*}

The same decoupling that makes this diagram work on one machine is what makes
the two-machine deployment in Section VII possible without changing a single
line of code. Every generator already reaches its target through an
environment variable (\texttt{TARGET\_URL}, \texttt{TARGET\_HOST}, or
\texttt{BROKER\_HOST}), never through a hardcoded address. Splitting the
system across two hosts is therefore just a matter of pointing those
variables at a real IP instead of a Docker service name, as described in
Section VII; no overlay network or Swarm setup is needed for two machines on
the same LAN.

The remaining five containers sit outside the request-handling pipeline shown
in Fig.~\ref{fig:architecture}: Mosquitto, already introduced above as the
target \texttt{gen-mqtt} and \texttt{target-tcpudp}'s silent counterpart, and
four \texttt{tshark}-based \textbf{analyzer} sidecars
(\texttt{analyzer-http2}, \texttt{analyzer-quic}, \texttt{analyzer-mqtt},
\texttt{analyzer-tcpudp}), one per target's network namespace
(\texttt{network\_mode: "service:target-X"}, so each sidecar sees exactly the
traffic its target sees, nothing more). Each runs a live capture filtered to
its own protocol's port, classifies packets purely by port number, and POSTs
a running total plus a one-second-bucketed history to the Metrics Collector
every second through the same \texttt{/analysis/update} endpoint the
generators already use for self-reported counters. This gives every number in
Section V two independent sources to cross-check against each other: what the
generators say they sent, and what \texttt{tshark} actually observed on the
wire.
